\begingroup
\raggedbottom
\setlength{\emergencystretch}{2em}
\section{双有理几何与正性}\label{src:B.2}
\begin{epigraph}
数学诸学科运用着由数字与符号组成的独特语言, 其精妙丝毫不逊于任何其他语言.
\epigraphauthor{Jorge Luis Borges}
\end{epigraph}
本节概述我们在发展模理论时所需的双有理几何基本事实. 在 $\overline{\mathcal M}_g$ 的稳定约化定理 (\bookref{6.5.1}) 中, \S\bookref{B.2.1-subsection} 的曲面双有理几何起着关键作用. 为了证明用于说明 $\overline M_g$ 射影的 Kollár 丰沛性判据 (\bookref{6.9.2}), 我们要用到 Nakai--Moishezon 丰沛性判据 (\bookref{B.2.28}) 以及 \S\bookref{B.2.4-subsection} 中数值有效向量丛的性质.

\subsection{曲面的双有理几何}\label{src:B.2.1-subsection}
对于整 Noether 概形 $X$, 一个\newterm{奇点消解}{resolution of singularities} 是从整正则概形 $X'$ 出发的固有双有理态射 $X'\to X$.
\begin{theorem}[消解的存在性]\label{src:B.2.1}
每个二维整 Noether 概形 $X$ 都有奇点消解.
\end{theorem}
\begin{proof}
特征 $0$ 的情形由 Zariski 证明 \cite{Zar39}, 特征 $p$ 的情形由 Abhyankar 证明 \cite{Abh56}, 混合特征的情形由 Lipman 证明 \cite{Lip78}. 另见 \cite[\S$2$]{Kol07} 及 \cite[Tag 0BGP]{SP}.
\end{proof}
\begin{theorem}[极小消解的存在性]\label{src:B.2.2}
设 $X$ 为二维整 Noether 概形. 存在奇点消解 $\pi:\widetilde X\to X$, 使得任意其他奇点消解 $Y\to X$ 都分解为 $Y\to\widetilde X\to X$. 此外, 对每条 $\pi$-例外曲线 $E$, 都有 $K_{\widetilde X}\cdot E\geq0$.
\end{theorem}
\begin{proof}
见 \cite[定理 $2.16$]{Kol07}.
\end{proof}
\begin{theorem}[嵌入消解的存在性]\label{src:B.2.3}
设 $X$ 为二维正则概形, $Y\subseteq X$ 为纯一维子概形. 假设对 $Y$ 的每个不可约分支 $Z$, 正规化 $\widetilde Z\to Z$ 都有限. 则存在在约化闭点处进行的有限次爆破
\[
X_n\longrightarrow\cdots\longrightarrow X_1\longrightarrow X,
\]
使得 $Y$ 的原像 $Y_n\subseteq X_n$ 是有效 Cartier 除子, 其支集为正规交叉除子, 即 $(Y_n)_{\mathrm{red}}$ 只有结点奇点.
\end{theorem}
\begin{proof}
见 \cite[定理 V.$3.9$]{Har77}, \cite[定理 $1.47$]{Kol07} 及 \cite[Tag 0BIC]{SP}.
\end{proof}
\begin{theorem}[双有理映射的分解]\label{src:B.2.4}
设 $X$ 和 $Y$ 均为二维正则整 Noether 概形. 每个固有双有理态射 $f:X\to Y$ 都是在约化闭点处进行的爆破的复合.
\end{theorem}
\begin{proof}
见 \cite[定理 V.$5.5$]{Har77}, \cite[定理 $2.13$]{Kol07} 及 \cite[Tag 0C5R]{SP}.
\end{proof}
\begin{theorem}[例外曲线的 Hodge 指数定理]\label{src:B.2.5}
设 $f:X\to Y$ 为二维 Noether 概形之间的射影且一般有限的态射, 其中 $X$ 正则, $Y$ 在某个域或 DVR 上拟射影. 设 $E_1,\ldots,E_n$ 为例外曲线. 则相交形式的矩阵 $(E_i\cdot E_j)$ 负定. 特别地, 对每个 $i$, 都有 $E_i^2<0$.
\end{theorem}
\begin{proof}
见 \cite[定理 $2.12$]{Kol07}.
\end{proof}
\begin{theorem}[Castelnuovo 收缩定理]\label{src:B.2.6}
设 $X$ 为二维正则概形, 且在域 $k$ 或剩余域为 $k$ 的 DVR $R$ 上射影, $E=\mathbb P^1_k\subseteq X$ 为满足 $E^2<0$ 的光滑有理曲线. 则存在到射影曲面的射影态射 $X\to Y$ 及一点 $y\in Y$, 使得 $f^{-1}(y)=E$, 且 $X\setminus E\to Y\setminus\{y\}$ 为同构. 若 $E^2=-1$, 则 $Y$ 光滑.
\end{theorem}
\begin{proof}
见 \cite[定理 V.$5.7$, 习题 V.$5.2$]{Har77} 及 \cite[定理 $2.14$, 注记 $2.15$]{Kol07}.
\end{proof}
可以证明, 在光滑射影曲面中反复收缩光滑有理 $-1$ 曲线的过程会终止 (见 \cite[定理 $5.8$]{Har77}). 因此, 有限次应用 Castelnuovo 可收缩性判据可得:
\begin{corollary}[极小模型的存在性]\label{src:B.2.7}
光滑曲面 $X$ 存在到光滑曲面的射影双有理态射 $X\to X_{\min}$, 使得任意到光滑曲面的射影双有理态射 $X_{\min}\to Y$ 都是同构. 特别地, $X_{\min}$ 不含光滑有理 $-1$ 曲线.
\end{corollary}

\subsection{正性}\label{src:B.2.2-subsection}
我们讨论线丛的正性性质, 其中一部分将在 \S\bookref{6.9.2-subsection} 推广到代数空间. 关于这些内容的优秀参考书是 \cite{Laz04a,Laz04b}.

\paragraph{丰沛线丛.}
概形 $X$ 上的线丛 $L$ 称为\newterm{丰沛}{ample}, 如果 $X$ 拟紧, 且对每个 $x\in X$, 都存在截面 $s\in\Gamma(X,L)$, 使得 $X_s=\{s\neq0\}$ 仿射且包含 $x$.
\begin{proposition}[丰沛性的刻画]\label{src:B.2.8}
对于 Noether 概形 $X$ 上的线丛 $L$, 下列条件等价:
\begin{enumerate}[(1)]
\item $L$ 丰沛;
\item 自然映射 $X\to\Proj\bigoplus_{d\geq0}\Gamma(X,L^{\otimes d})$ 良定义且为开浸入;
\item 对每个凝聚层 $F$, 张量积 $F\otimes L^{\otimes m}$ 在 $m\gg0$ 时无基点.
\end{enumerate}
若进一步假设 $X$ 在 Noether 环 $R$ 上固有, 则以上条件还等价于:
\begin{enumerate}[(1)]
\setcounter{enumi}{3}
\item 对某个 $m>0$, $L^{\otimes m}$ 非常丰沛, 即定义了到射影空间的闭嵌入 $|L^{\otimes m}|:X\hookrightarrow\mathbb P^N_R$;
\item 对 $X$ 上每个凝聚层 $F$, 当 $m\gg0$ 且 $i>0$ 时, $\mathrm H^i(X,F\otimes L^{\otimes m})=0$.
\end{enumerate}
\end{proposition}
\begin{proof}
见 \cite[\S II.$7$ 及 III.$5.3$]{Har77}, \cite[II.$4.5$ 及 III.$2.6$]{EGA}, 或 \cite[Tags 01PR 及 0B5U]{SP}.
\end{proof}
\begin{proposition}[丰沛性的性质]\label{src:B.2.9}
设 $X$ 为域 $k$ 上的固有概形, $L$ 为 $X$ 上的线丛.
\begin{enumerate}[(1)]
\item 若 $f:X'\to X$ 为有限满射态射, 则 $L$ 丰沛当且仅当 $f^*L$ 丰沛.
\item 对于域扩张 $k\to k'$, $L$ 在 $X$ 上丰沛当且仅当 $L_{k'}$ 在 $X_{k'}$ 上丰沛.
\end{enumerate}
\end{proposition}
\begin{proof}
两者都由丰沛性的上同调刻画推出. 对于 (1), 见 \cite[习题 III.$5.7$]{Har77}, \cite[III.$2.6.2$]{EGA} 及 \cite[Tag 0B5V]{SP}. (2) 也直接由丰沛性的 fpqc 下降 (\bookref{B.2.12}) 推出.
\end{proof}
(1) 蕴含 $X$ 上的线丛 $L$ 丰沛当且仅当它在每个不可约分支 $X_i$ 的约化子概形上的限制 $L|_{(X_i)_{\mathrm{red}}}$ 丰沛.
\begin{proposition}[丰沛性的开放性]\label{src:B.2.10}
设 $f:X\to S$ 为概形之间的固有有限呈示态射, $L$ 为 $X$ 上的线丛. 若对某个 $s\in S$, $L$ 在纤维 $X_s$ 上的限制 $L_s$ 丰沛 (相应地, 非常丰沛且对 $i>0$ 有 $\mathrm H^i(X_s,L_s)=0$), 则存在 $s$ 的开邻域 $U\subseteq S$, 使得 $L_U$ 在 $X_U$ 上相对于 $U$ 丰沛 (相应地, 相对非常丰沛). 特别地, 对所有 $u\in U$, $L_u$ 在 $X_u$ 上丰沛 (相应地, 非常丰沛).
\end{proposition}
\begin{proof}
我们在 $X\to S$ 平坦的条件下给出证明 (这已足以应付许多应用). 若 $L_s$ 在 $X_s$ 上丰沛, 则对 $n\gg0$, $L_s^{\otimes n}$ 非常丰沛且对 $i>0$ 有 $\mathrm H^i(X_s,L_s^{\otimes n})=0$. 因此只须处理非常丰沛的情形. 由上同调与基变换 (\bookref{A.6.8}), 将 $S$ 换为 $s$ 的一个开邻域后, $f_*L$ 为向量丛, 且对 $t\in S$, 比较映射 $f_*L\otimes\kappa(t)\to\mathrm H^i(X_t,L_t)$ 为同构. 进一步换为仿射开邻域后, 可使 $\mathrm H^0(X,L)$ 由截面 $t_0,\ldots,t_n$ 自由生成, 它们在 $\mathrm H^0(X_s,L_s)$ 中的限制构成一组基. 零点集 $V:=V(t_0,\ldots,t_n)\subseteq X$ 闭且与 $X_s$ 不交. 将 $S$ 换为包含在 $S\setminus f(V)$ 中的一个 $s$ 的仿射开邻域, 可假设截面 $t_i$ 生成 $L$, 并定义 $S$ 上的态射 $g:X\to\mathbb P^n_S$, 其限制为闭浸入 $g_s:X_s\hookrightarrow\mathbb P^n_{\kappa(s)}$.

由纤维维数的上半连续性, 满足 $\dim g^{-1}(z)>0$ 的点 $z$ 构成闭集 $Z\subseteq\mathbb P^n_S$. 由于 $Z$ 与 $\mathbb P^n_{\kappa(s)}$ 不交, 可进一步缩小 $S$, 使 $g:X\to\mathbb P^n_S$ 拟有限, 从而有限 (因为 $g$ 固有). $\mathcal O_{\mathbb P^n_S}\to g_*\mathcal O_X$ 的余核凝聚, 其支集为 $\mathbb P^n$ 的闭子概形, 且与 $\mathbb P^n_{\kappa(s)}$ 不交. 进一步缩小 $S$ 后, 可使 $g:X\to\mathbb P^n_S$ 为闭浸入, 因而 $L=g^*\mathcal O_{\mathbb P^n_S}(1)$ 非常丰沛. 另见 \cite[定理 $1.2.17$, 定理 $1.7.8$]{Laz04a}, \cite[III$_1$.$4.7.1$, IV$_3$.$9.6.4$]{EGA}, \cite[命题 $1.41$]{KM98} 及 \cite[Tag 0D3D]{SP}.
\end{proof}
\begin{example}[由 Brian Nugent 提出]\label{src:B.2.11}
非常丰沛性并非开条件. 若 $C$ 为亏格 $3$ 的非超椭圆曲线, 则 $K_C$ 非常丰沛, 且可写成 $K_C=\mathcal O_C(p_1+\cdots+p_4)$. 考虑常值族 $C\times C\to C$, 配备常值截面 $p_1,p_2,p_3$ 和对角截面 $\Delta$. 则 $\mathcal O_{C\times C}(p_1+p_2+p_3+\Delta)$ 在 $p_4\in C$ 上的纤维非常丰沛, 而在 $p_4$ 附近的点 $s\in C$ 上的纤维并不非常丰沛.
\end{example}
我们还需要丰沛性是靶上的 fpqc 局部性质这一事实.
\begin{proposition}[丰沛性的 fpqc 下降]\label{src:B.2.12}
设 $f:X\to S$ 为概形态射, $L$ 为 $X$ 上的线丛. 若 $S'\to S$ 为概形之间的 fpqc 态射, 则 $L$ 相对于 $S$ 丰沛当且仅当它到 $X\times_S S'$ 的拉回相对于 $S'$ 丰沛.
\end{proposition}
\begin{proof}
我们知道相对丰沛性在基变换下稳定. 若到 $X':=X\times_S S'$ 的拉回 $L'$ 相对丰沛, 则 $X'\to S'$ 拟紧且拟分离. 由于它们都是靶上的 fpqc 局部性质 (\bookref{3.1.26}), $f:X\to S$ 也拟紧且拟分离. 因此可构造拟凝聚分次 $\mathcal O_S$-代数 $\mathcal A=\bigoplus_{d\geq0}f_*L^{\otimes d}$ 及相应态射 $X\to\Proj_S\mathcal A$. 我们使用如下事实: $L$ 相对丰沛当且仅当 $X\to\Proj_S\mathcal A$ 良定义且为开浸入. 此态射的基变换为 $X'\to\Proj_{S'}\mathcal A'$, 其中 $\mathcal A'=\bigoplus_{d\geq0}f'_*L'^{\otimes d}$. 由于开浸入是靶上的 fpqc 局部性质 (\bookref{3.1.26}), 结论成立. 另见 \cite[IV$_2$.$2.7.2$]{EGA} 及 \cite[Tag 0D2P]{SP}.
\end{proof}

\paragraph{数值有效线丛.}
域 $k$ 上固有概形 $X$ 的线丛 $L$ 称为\newterm{数值有效}{nef} (或 numerically effective), 如果对每条整闭曲线 $C\subseteq X$, 都有
\[
\int_C c_1(L)\geq0.
\]
这里 $\int_C c_1(L)$ 表示的量与 $c_1(L)\cdot C$, $L\cdot C$ 或 $\deg L|_C$ 相同. 若对每条整闭曲线都有 $c_1(L)\cdot C>0$, 则称 $L$ \newterm{严格数值有效}{strictly nef}.
\begin{theorem}[Kleiman 定理]\label{src:B.2.13}
若 $L$ 为域 $k$ 上固有概形 $X$ 的线丛, 则 $L$ 数值有效当且仅当对每个 $k$ 维整子概形 $Z\subseteq X$, 都有
\[
\int_Z c_1(L)^k\geq0.
\]
\end{theorem}
\begin{proof}
见 \cite[定理 $1.4.9$]{Laz04a}, \cite[定理 $2.17$]{Kol96}, 或原始文献 \cite{Kle66}.
\end{proof}
采用加法记号表示线丛往往很方便, 此时 $mL+H$ 对应于 $L^{\otimes m}\otimes H$.
\begin{corollary}[数值有效性的刻画]\label{src:B.2.14}
设 $X$ 为域 $k$ 上的射影概形, $H$ 为丰沛线丛. $X$ 上的线丛 $L$ 数值有效当且仅当 $mL+H$ 在 $m\gg0$ 时丰沛.
\end{corollary}
\begin{proof}
见 \cite[推论 $1.4.10$]{Laz04a}.
\end{proof}
\begin{proposition}[数值有效性的性质]\label{src:B.2.15}
设 $X$ 为域 $k$ 上的固有概形, $L$ 为 $X$ 上的线丛.
\begin{enumerate}[(1)]
\item 若 $f:X'\to X$ 为满射固有态射, 则 $L$ 数值有效当且仅当 $f^*L$ 数值有效.
\item 对于域扩张 $k\to k'$, $L$ 在 $X$ 上数值有效当且仅当 $L_{k'}$ 在 $X_{k'}$ 上数值有效.
\end{enumerate}
\end{proposition}
\begin{proof}
对于 (1), 若 $C'\subseteq X'$ 为整曲线, $d$ 为诱导映射 $C'\to f(C')$ 的次数, 则由投影公式有
\begin{equation}\label{src:B.2.16}
f^*L\cdot C'=d\bigl(L\cdot f(C')\bigr).
\end{equation}
这给出 ($\Rightarrow$). 反之, 若 $C\subseteq X$ 为整曲线, 可选取整曲线 $C'\subseteq X'$ 使 $C=f(C')$, 从而 \bookeqref{B.2.16} 也给出 ($\Leftarrow$). 对于 (2), 由 Chow 引理及 (1), 可假设 $X$ 射影. 此时数值有效性的刻画 (\bookref{B.2.14}) 将问题化为相应的丰沛性结论 (命题 \bookref{B.2.9}(2)).
\end{proof}
\begin{proposition}[数值有效性在一般化下稳定]\label{src:B.2.17}
设 $X$ 为 DVR $R$ 上的固有平坦概形, $L$ 为 $X$ 上的线丛. 若 $L$ 在中心纤维 $X_0$ 上的限制 $L_0$ 数值有效, 则它在一般纤维 $X_\eta$ 上的限制 $L_\eta$ 也数值有效.
\end{proposition}
\begin{proof}
由 Chow 引理及命题 \bookref{B.2.15}(1), 可假设 $X$ 射影, 并有丰沛线丛 $H$. 由数值有效性的刻画 (\bookref{B.2.14}), $mL_0+H_0$ 在 $m\gg0$ 时丰沛. 由丰沛性的开放性 (\bookref{B.2.10}), $mL+H$ 丰沛, 因而 $mL_\eta+H_\eta$ 也丰沛. 再次应用数值有效性的刻画, 即知 $L_\eta$ 数值有效.
\end{proof}
\begin{remark}\label{src:B.2.18}
对于代数簇之间的满射固有态射 $X\to S$ 及 $X$ 上的线丛 $L$, 若 $L$ 在 $s\in S$ 上的纤维 $L_s$ 数值有效, 则存在不包含 $s$ 的真子概形的可数并 $B\subseteq S$, 使得对每个 $t\in S\setminus B$, $L_t$ 都数值有效 \cite[命题 $1.4.14$]{Laz04a}. 数值有效性一般并非开条件; 见 \cite[例 $5.3$]{Lan13} 及 \cite[定理 $1.2$]{Les14}.
\end{remark}
\begin{remark}[丰沛锥与数值有效锥]\label{src:B.2.19}
丰沛线丛和数值有效线丛分别生成 $N^1(X)_{\mathbb R}$ 中的锥 $\operatorname{Amp}(X)$ 和 $\operatorname{Nef}(X)$, 称为\newterm{丰沛锥}{ample cone} 和\newterm{数值有效锥}{nef cone}. 对于射影代数簇, Kleiman 定理 (\bookref{B.2.13}) 蕴含数值有效锥是丰沛锥的闭包, 而丰沛锥是数值有效锥的内部; 见 \cite[定理 $1.4.23$]{Laz04a}.
\end{remark}

\paragraph{有效, 无基点与半丰沛线丛.}
对于域 $k$ 上固有概形 $X$ 的线丛 $L$, 有下列概念:
\begin{itemize}
\item 若 $\Gamma(X,L)\neq0$, 则称 $L$ \textit{有效} (\textup{effective})\index{effective line bundle@有效线丛\quad effective line bundle|hyperpage}.
\item 若对每个 $x\in X$, 都存在 $s\in\Gamma(X,L)$ 使 $s(x)\neq0$, 则称 $L$ \newterm{无基点}{base point free} (或\newterm{整体生成}{globally generated}); 等价地, 完全线性系 $|L|$ 定义态射 $X\to\mathbb P(\mathrm H^0(X,L))$.
\item 若对某个 $m>0$, $L^{\otimes m}$ 无基点, 则称 $L$ \newterm{半丰沛}{semiample}.
\end{itemize}
半丰沛线丛 $L$ 必然数值有效. 事实上, 若对某个 $m>0$, $L^{\otimes m}$ 定义态射 $f:X\to\mathbb P^N$, 且 $f^*\mathcal O(1)\simeq c_1(L^{\otimes m})$, 则投影公式给出 $C\cdot L^{\otimes m}=f(C)\cdot c_1(\mathcal O(1))\geq0$. 因此有如下蕴含关系:
\[
\begin{tikzcd}[column sep=2.6em]
\text{无基点}\arrow[r,Rightarrow]&\text{半丰沛}\arrow[r,Rightarrow]&\text{数值有效}.
\end{tikzcd}
\]

\paragraph{大线丛.}
域 $k$ 上固有概形 $X$ 的线丛 $L$ 称为\newterm{大}{big}, 如果存在常数 $C>0$, 使得对 $m\gg0$ 有 $h^0(X,L^{\otimes m})>C\cdot m^{\dim(X)}$.
\begin{proposition}[Kodaira 引理]\label{src:B.2.20}
设 $X$ 为域 $k$ 上的射影概形, $L$ 为 $X$ 上的大线丛. 若 $E$ 为有效线丛, 则对充分可除的 $m$, $mL-E$ 有效.
\end{proposition}
\begin{proof}
见 \cite[命题 $2.2.6$]{Laz04a}.
\end{proof}
\begin{proposition}[大性的刻画]\label{src:B.2.21}
对于域 $k$ 上的射影整概形 $X$ 及其上的线丛 $L$, 下列条件等价:
\begin{enumerate}[(1)]
\item $L$ 大;
\item 对某个 $m>0$, $|L^{\otimes m}|$ 定义有理映射 $X\dashrightarrow\mathbb P(\mathrm H^0(X,L^{\otimes m}))$, 且它到其像是双有理的;
\item 对 $X$ 上每个丰沛除子 $A$, 都存在正整数 $m>0$ 和 $X$ 上的有效除子 $N$, 使得 $mL=A+N$ (线性等价);
\item 存在 $X$ 上的丰沛除子 $A$, 正整数 $m>0$ 及 $X$ 上的有效除子 $N$, 使得 $mL\equiv A+N$ (数值等价);
\item 存在整数 $m>0$ 及整体截面 $s\in\mathrm H^0(X,L^{\otimes m})$, 使得 $X_s$ 非空且拟仿射 (相应地, 仿射).
\end{enumerate}
\end{proposition}
\begin{proof}
见 \cite[推论 $2.2.7$]{Laz04a} 及 \cite[引理 $2.8$--$2.9$]{IO25}.
\end{proof}
由此可知, 在允许缩放 (即取正张量幂) 的意义下, 大线丛就是丰沛线丛与有效线丛之和. 这蕴含大线丛与有效线丛之和仍然大. 归纳如下:
\[
\begin{tikzcd}[column sep=5em,row sep=1.8em]
\text{大}\arrow[r,Leftrightarrow,"\text{允许缩放}"]&\text{丰沛}+\text{有效}\\
\text{大}+\text{有效}\arrow[r,Rightarrow]&\text{大}.
\end{tikzcd}
\]
\begin{theorem}[渐近 Riemann--Roch 定理]\label{src:B.2.22}
设 $X$ 为域 $k$ 上的 $n$ 维固有概形, $L$ 为 $X$ 上的线丛. 则 Euler 示性数
\[
\chi(X,L^{\otimes m})=\frac{(c_1(L)^n)}{n!}m^n+O(m^{n-1})
\]
是关于 $m$ 的次数不超过 $n$ 的多项式. 若进一步假设 $L$ 数值有效, 则 $h^i(X,L^{\otimes m})=O(m^{n-1})$, 且
\[
h^0(X,L^{\otimes m})=\frac{(c_1(L)^n)}{n!}m^n+O(m^{n-1}).
\]
\end{theorem}
\begin{proof}
见 \cite[推论 $1.4.41$]{Laz04a} (射影情形), \cite[定理 VI.$2.14$--$15$]{Kol96} 及 \cite[Tag 0BJ8]{SP}.
\end{proof}
承认 $\chi(X,L^{\otimes m})$ 是多项式后, 可以将相交数 $c_1(L)^n$ 定义为其规范化首项系数. 也可以用其他几种方式定义它, 参见 \cite[定理 VI.$2$]{Kol96}. 渐近 Riemann--Roch 定理给出了数值有效线丛的大性的一个有用刻画, 例如它蕴含丰沛线丛是大的.
\begin{corollary}[大性的刻画 II]\label{src:B.2.23}
设 $X$ 为域 $k$ 上的 $n$ 维固有概形. $X$ 上的数值有效线丛 $L$ 大当且仅当 $c_1(L)^n>0$.
\end{corollary}
\begin{proposition}[大性的性质]\label{src:B.2.24}
设 $X$ 为域 $k$ 上的整固有概形, $L$ 为 $X$ 上的线丛.
\begin{enumerate}[(1)]
\item 设 $f:X'\to X$ 为一般拟有限的固有概形态射. 则 $L$ 大当且仅当 $f^*L$ 大.
\item 对于域扩张 $k\to k'$, $L$ 在 $X$ 上大当且仅当 $L_{k'}$ 在 $X_{k'}$ 上大.
\end{enumerate}
\end{proposition}
\begin{proof}
对于 (1), 投影公式给出
\begin{align*}
\mathrm H^0(X',f^*L^{\otimes m})
&=\mathrm H^0(X,f_*f^*L^{\otimes m})\\
&=\mathrm H^0(X,L^{\otimes m}\otimes f_*\mathcal O_{X'}).
\end{align*}
由于 $\mathcal O_X\hookrightarrow f_*\mathcal O_{X'}$ 单射, 有包含 $\mathrm H^0(X,L^{\otimes m})\hookrightarrow\mathrm H^0(X,L^{\otimes m}\otimes f_*\mathcal O_{X'})$. 因为 $\dim X'=\dim X$, $L$ 大蕴含 $f^*L$ 大. 注意, 若 $X'\to X$ 双有理且 $X$ 正规, 则 $\mathcal O_X=f_*\mathcal O_{X'}$, 且 $\mathrm H^0(X',f^*L^{\otimes m})=\mathrm H^0(X,L^{\otimes m})$, 从而得到逆命题. 一般情形下, 相交数的投影公式给出 $c_1(f^*L)^{\dim X'}=\deg(f)c_1(L)^{\dim X}$. 由于 $L$ 数值有效当且仅当 $f^*L$ 数值有效 (命题 \bookref{B.2.15}), 推论 \bookref{B.2.23} 蕴含 $L$ 大当且仅当 $f^*L$ 大. (2) 由等式 $\mathrm H^0(X',f^*L^{\otimes m})=\mathrm H^0(X,L^{\otimes m})\otimes_k k'$ 推出.
\end{proof}
\begin{remark}[大锥与伪有效锥]\label{src:B.2.25}
大除子与有效除子在 $N^1(X)_{\mathbb R}$ 中分别生成\newterm{大锥}{big cone} $\operatorname{Big}(X)$ 和\newterm{有效锥}{effective cone} $\operatorname{Eff}(X)$, 其闭包 $\overline{\operatorname{Eff}}(X)$ 称为\newterm{伪有效锥}{pseudo-effective cone}. 大锥 $\operatorname{Big}(X)$ 包含于 $\overline{\operatorname{Eff}}(X)$ 的内部, 且 $\overline{\operatorname{Eff}}(X)=\overline{\operatorname{Big}}(X)$ \cite[定理 $2.2.6$]{Laz04a}.
\end{remark}

\subsection{丰沛性判据}\label{src:B.2.3-subsection}
我们回顾检验固有概形上线丛丰沛性的技巧. 首先要记住的策略是: 半丰沛且严格数值有效 $\Rightarrow$ 丰沛.
\begin{lemma}\label{src:B.2.26}
在域 $k$ 上的固有概形上, 线丛 $L$ 丰沛当且仅当它严格数值有效且半丰沛.
\end{lemma}
\begin{proof}
对于非平凡的方向, 存在 $m>0$ 使 $L^{\otimes m}$ 定义态射 $f:X\to\mathbb P^N$, 且它不收缩任何曲线. 因此 $f:X\to\mathbb P^N$ 固有且拟有限, 从而有限. 故 $L^{\otimes m}=f^*\mathcal O(1)$ 丰沛.
\end{proof}
\begin{remark}\label{src:B.2.27}
实践中, 半丰沛性条件可能很难检验. 双有理几何中有一些强有力的无基点定理, 有时可用它们将半丰沛性化为大性与数值有效性. 例如, Kawamata 无基点定理断言: 若 $(X,\Delta)$ 为固有 klt 对, $\Delta$ 有效, 且 $D$ 为数值有效 Cartier 除子, 使得对某个 $a>0$, $aD-K_X-\Delta$ 数值有效且大, 则 $D$ 半丰沛 \cite[定理 $3.3$]{KM98}. 可将此结果与\newterm{丰度猜想}{Abundance Conjecture} 对照: 若 $(X,\Delta)$ 为固有对数典范对, $\Delta$ 有效, 则 $K_X+\Delta$ 数值有效蕴含其半丰沛 \cite[猜想 $3.12$]{KM98}.

另一种方法来自 Zariski 和 Wilson 的经典结果: 若 $X$ 为正规射影代数簇, $D$ 为数值有效大除子, 则 $D$ 半丰沛当且仅当其分次截面环 $\bigoplus_n\Gamma(X,\mathcal O_X(nD))$ 有限生成; 见 \cite[定理 $2.3.15$]{Laz04a}. 虽然有时可应用 \cite{BCHM10} 检验有限生成性, 这一结果已经预先假定了 $X$ 的射影性. 尽管如此, 它有时仍然有用. 在正特征下, Keel 定理 \cite{Kee99} 提供了另一种方法: 在特征 $p$ 的域 $k$ 上的射影代数簇 $X$ 上, 数值有效线丛 $L$ 半丰沛当且仅当它在例外集 $E$ 上的限制半丰沛, 其中例外集 $E$ 定义为满足 $L^{\dim Z}\cdot Z=0$ 的不可约子簇 $Z\subseteq X$ 的并.
\end{remark}

\paragraph{丰沛性的数值判据.}
Nakai--Moishezon 丰沛性判据\footnote{它也称为 Nakai 判据或 Nakai--Moishezon--Kleiman 判据. 历史介绍与进一步的参考文献见 \cite[\S$1.2$.B]{Laz04a}.} 提供了证明射影性的便捷方法. 我们将在定理 \bookref{6.9.10} 中将此判据推广到代数空间, 因为它用于 Kollár 判据 (\bookref{6.9.2}) 的证明, 而后者又用于证明 $\overline M_g$ 的射影性.
\begin{theorem}[Nakai--Moishezon 丰沛性判据]\label{src:B.2.28}
设 $X$ 为代数闭域 $k$ 上的固有概形, $L$ 为 $X$ 上的线丛. 下列条件等价:
\begin{enumerate}[(1)]
\item $L$ 丰沛;
\item 对每个整闭子概形 $Z\subseteq X$, 都有 $c_1(L)^{\dim Z}\cdot Z>0$;
\item $L$ 数值有效, 且对每个整闭子概形 $Z\subseteq X$, $L|_Z$ 都大;
\item $L$ 严格数值有效, 且对每个整闭子概形 $Z\subseteq X$, 存在 $m>0$ 使 $L|_Z^{\otimes m}$ 有效.
\end{enumerate}
\end{theorem}
\begin{proof}
数值有效线丛 $L$ 大当且仅当最高相交数为正 (推论 \bookref{B.2.23}). 这给出 (2) 与 (3) 的等价性. 因而有 $(1)\Rightarrow(2)\Leftrightarrow(3)\Rightarrow(4)$. 对于 $(4)\Rightarrow(1)$, 由于 $L$ 严格数值有效, 由引理 \bookref{B.2.26} 只须验证 $L$ 半丰沛. 对某个除子 $D$, 写成 $L=\mathcal O_X(D)$. 因为 $D$ 在 $X$ 上大, 某个正整数倍 $mD$ 有效. 将 $D$ 换成 $mD$ 后, 取非零截面 $s\in\mathrm H^0(X,\mathcal O_X(D))$. 则 $\mathcal O_X(D)$ 在 $X\setminus D$ 上无基点. 我们将证明对 $m\gg0$, $\mathcal O_X(mD)$ 在 $D$ 上也无基点. 对 $\dim X$ 归纳, 可假设 $\mathcal O_X(D)|_D$ 丰沛. 考虑正合列
\[
0\longrightarrow\mathcal O_X((m-1)D)\longrightarrow\mathcal O_X(mD)
\longrightarrow\mathcal O_X(mD)|_D\longrightarrow0.
\]
\setlength{\emergencystretch}{4em}
对 $m\gg0$, $\mathcal O_X(mD)|_D$ 无基点且 $\mathrm H^1(X,\mathcal O_X(mD)|_D)=0$. 因此 $\mathrm H^1(X,\mathcal O_X((m-1)D))\twoheadrightarrow\mathrm H^1(X,\mathcal O_X(mD))$ 满射; 由于每个向量空间均有限维, 这些满射在 $m\gg0$ 时最终成为同构. 因而对 $m\gg0$, $\mathrm H^0(X,\mathcal O_X(mD))\to\mathrm H^0(D,\mathcal O_X(mD)|_D)$ 满射, $\mathcal O_X(mD)$ 在 $D$ 上无基点. 另见 \cite[定理 V.$1.10$]{Har77} (曲面情形), \cite[定理 $1.2.23$]{Laz04a} (射影情形), \cite[定理 VI.$2.18$]{Kol96}, \cite[定理 $1.37$]{KM98} 及 \cite[\S III.$1$]{Kle66}.
\end{proof}
\setlength{\emergencystretch}{2em}
虽然本书不会应用下面的判据, 但它们在其他场合往往很有用.
\begin{theorem}[Kleiman 判据]\label{src:B.2.29}
若 $X$ 为代数闭域 $k$ 上的射影概形, 或为 $\mathbb Q$-阶乘 (例如光滑) 固有概形, 则 $X$ 上的线丛 $L$ 丰沛当且仅当对所有 $C\in\overline{\operatorname{Eff}}(X)$, 都有 $c_1(L)\cdot C>0$.
\end{theorem}
\begin{proof}
见 \cite[\S III.$1$]{Kle66}, \cite[定理 VI.$2.19$]{Kol96}, \cite[定理 $1.18$]{KM98} 及 \cite[定理 $1.4.23$]{Laz04a}.
\end{proof}
注意, 仅对整曲线 $C\subseteq X$ 检验 $c_1(L)\cdot C>0$ 并不足够; 必须对有效曲线锥的闭包 $\overline{\operatorname{Eff}}(X)$ 中的曲线类检验这一条件. Mumford 给出的反例见 \cite[第 $50$--$56$ 页]{Har70}.
\begin{theorem}[Seshadri 判据]\label{src:B.2.30}
若 $X$ 为代数闭域 $k$ 上的固有概形, 则 $X$ 上的线丛 $L$ 丰沛当且仅当存在 $\epsilon>0$, 使得对每个点 $x\in X$ 和每条整曲线 $C\subseteq X$, 都有 $c_1(L)\cdot C>\epsilon\operatorname{mult}_x(C)$, 其中 $\operatorname{mult}_x(C)$ 表示 $C$ 在 $x$ 处的重数.
\end{theorem}
\begin{proof}
见 \cite[定理 $1.4.13$]{Laz04a} 及 \cite[定理 $2.18$]{Kol96}.
\end{proof}

\subsection{数值有效向量丛}\label{src:B.2.4-subsection}
在 Kollár 丰沛性判据 (\bookref{6.9.2}) 中, 向量丛的数值有效性起着关键作用.
\begin{definition}\label{src:B.2.31}
概形 $X$ 上的向量丛 $E$ 称为\newterm{数值有效向量丛}{nef vector bundle} (或\newterm{半正向量丛}{semipositive vector bundle}), 如果 $\mathcal O_{\mathbb P(E)}(1)$ 在 $\mathbb P(E)$ 上数值有效.
\end{definition}
还有相关的丰沛向量丛概念, 定义为要求 $\mathcal O_{\mathbb P(E)}(1)$ 在 $\mathbb P(E)$ 上丰沛; 本书不需要这一概念, 见 \cite{Har66a} 及 \cite[\S$6$]{Laz04b}.
\begin{proposition}[向量丛数值有效性的刻画]\label{src:B.2.32}
设 $E$ 为代数闭域 $k$ 上固有概形 $X$ 的向量丛. 下列条件等价:
\begin{enumerate}[(1)]
\item $E$ 数值有效;
\item 对每个从光滑固有曲线出发的映射 $f:C\to X$, 每个商线丛 $f^*E\twoheadrightarrow L$ 都有非负次数;
\item 对每个从光滑固有曲线出发的映射 $f:C\to X$, 每个商向量丛 $f^*E\twoheadrightarrow W$ 都有非负次数.
\end{enumerate}
\end{proposition}
\begin{proof}
见 \cite[第 $437$ 页]{Bar71}, \cite[命题 $6.1.18$]{Laz04b} 及 \cite[定义-命题 $3.3$]{Kol90}.
\end{proof}
\begin{proposition}[向量丛数值有效性的性质]\label{src:B.2.33}
设 $X$ 为域 $k$ 上的固有概形, $E$ 为 $X$ 上的向量丛.
\begin{enumerate}[(1)]
\item 若 $f:X'\to X$ 为满射固有态射, 则 $E$ 数值有效当且仅当 $f^*E$ 数值有效.
\item 对于域扩张 $k\to k'$, $E$ 在 $X$ 上数值有效当且仅当 $E_{k'}$ 在 $X_{k'}$ 上数值有效.
\item 数值有效向量丛的商, 扩张及张量积均数值有效. 若 $E$ 数值有效, 则对 $k\geq0$, $\bigwedge^kE$, $\operatorname{Sym}^kE$ 及 $\operatorname{Sym}^k(E^\vee)^\vee$ 也数值有效.
\end{enumerate}
\end{proposition}
\begin{proof}
(1) 和 (2) 由数值有效线丛的相应性质 (命题 \bookref{B.2.15}) 推出. (3) 需要一些工作. 由 (2), 可假设 $k$ 代数闭. 由向量丛数值有效性的刻画 (\bookref{B.2.32}), 只须假设 $X$ 为亏格 $g$ 的光滑曲线. 由此刻画可清楚看出, 数值有效丛的商仍数值有效; 数值有效丛的扩张数值有效也不难证明.

在证明其余部分前, 我们断言: 若 $E$ 为 $X$ 上的数值有效向量丛, $L$ 为满足 $\deg L\geq2g-1$ (相应地, $\deg L\geq2g$) 的线丛, 则 $\mathrm H^1(X,E\otimes L)=0$ (相应地, $E\otimes L$ 整体生成). 由 Serre 对偶 (\bookref{6.1.3}), 有
\[
\mathrm H^1(X,E\otimes L)=\operatorname{Hom}_{\mathcal O_X}(E\otimes L,\omega_X).
\]
若 $E\otimes L\to\omega_X$ 为非零映射, 则其像 $I\subseteq\omega_X$ 为线丛, 且 $\deg I\leq\deg\omega_X=2g-2$. 若 $\deg L\geq2g-1$, 则诱导商 $E\twoheadrightarrow I\otimes L^\vee$ 的次数为负, 与 $E$ 数值有效矛盾, 故 $\mathrm H^1(X,E\otimes L)=0$. 若 $\deg L\geq2g$, 则对每个 $p\in X(\mathrm{base})$, 有 $\mathrm H^1(X,E\otimes L(-p))=0$, 且 $\mathrm H^0(X,E\otimes L)\to\mathrm H^0(X,E\otimes L\otimes\kappa(p))$ 满射.

设 $E$ 和 $F$ 为 $X$ 上的数值有效丛, 秩分别为 $e$ 和 $f$. 取次数 $d\geq2g$ 的线丛 $L$. 由于 $E\otimes L$ 整体生成, 可选取 $E\otimes L$ 的 $e$ 个整体截面, 使它们在一般点 $\xi\in X$ 上的限制构成 $E\otimes L\otimes\kappa(\xi)$ 的一组基. 这给出一般满射的映射 $(L^\vee)^{\oplus e}\to E$. 同样, 有一般满射的映射 $(L^\vee)^{\oplus f}\to F$. 取这些映射的张量积, 得到一般满射的映射
\[
(L^\vee)^{\oplus ef}\longrightarrow E\otimes F.
\]
若 $Q$ 为 $E\otimes F$ 的商线丛, 则 $(L^\vee)^{\oplus ef}$ 在 $Q$ 中的像 $Q'\subseteq Q$ 满足 $\deg Q\geq\deg Q'\geq-d$. 因此, 任意秩分别为 $e$ 和 $f$ 的数值有效丛 $E$ 和 $F$, 其张量积 $E\otimes F$ 的每个商线丛 $Q$ 都满足 $\deg Q\geq-d$.

设 $\operatorname{char}(k)=p$. 假设存在商线丛 $E\otimes F\twoheadrightarrow Q$, 满足 $\deg Q<0$. 记 $\operatorname{Fr}^N:X\to X$ 为绝对 Frobenius 的 $N$ 次幂, 则商
\[
(\operatorname{Fr}^N)^*E\otimes(\operatorname{Fr}^N)^*F
\twoheadrightarrow(\operatorname{Fr}^N)^*Q
\]
为次数 $Np\deg Q$ 的商线丛. 由命题 \bookref{B.2.33}(1), $(\operatorname{Fr}^N)^*E$ 和 $(\operatorname{Fr}^N)^*F$ 数值有效. 取 $N$ 使 $Np\deg Q<-d$, 就得到次数小于 $-d$ 的商线丛, 与上述事实矛盾.

设 $\operatorname{char}(k)=0$. 由于 $X$ 在 $k$ 上有限型, 其定义方程只涉及 $k$ 中有限多个系数. 因此存在有限生成的 $\mathbb Z$-子代数 $A\subseteq k$, $A$ 上的有限型概形 $\widetilde X$, 以及 Cartesian 图
\[
\begin{tikzcd}[column sep=3em,row sep=2em]
X\arrow[r]\arrow[d]\arrow[dr,phantom,"\square" description]&\widetilde X\arrow[d]\\
\Spec k\arrow[r]&\Spec A.
\end{tikzcd}
\]
由极限下态射性质的下降 (\bookref{B.3.7}), 还可使 $\widetilde X\to\Spec A$ 为光滑曲线族. 最后, 沿映射 $\Spec R\to\Spec A$ 限制, 可假设 $A$ 为 DVR, 且其闭点与一般点的特征分别为 $p$ 和 $0$. 因而, 利用向量丛的数值有效性在一般化下稳定 (\bookref{B.2.34}), 可归结为正特征情形. $\bigwedge^kE$, $\operatorname{Sym}^kE$ 及 $\operatorname{Sym}^k(E^\vee)^\vee$ 的数值有效性由相同论证推出. 实际上, 这一论证给出更一般的结论: 若 $E_1,\ldots,E_m$ 为秩分别为 $r_1,\ldots,r_m$ 的数值有效向量丛, 表示
\[
\rho:\mathrm{GL}_{r_1}\times\cdots\times\mathrm{GL}_{r_m}\longrightarrow\mathrm{GL}_N
\]
半正, 即可延拓为映射
\[
\operatorname{Mat}_{r_1,r_1}\times\cdots\times\operatorname{Mat}_{r_m,r_m}
\longrightarrow\operatorname{Mat}_{N,N},
\]
则诱导向量丛 $\rho(E_1,\ldots,E_m)$ 也数值有效. 另见 \cite[命题 $3.5$]{Kol90}, \cite[定理 $3.3$]{Bar71} 及 \cite[定理 $5.2$]{Har66a}.
\end{proof}
\begin{proposition}[向量丛的数值有效性在一般化下稳定]\label{src:B.2.34}
设 $X$ 为 DVR $R$ 上的固有平坦概形, $E$ 为 $X$ 上的向量丛. 若 $E$ 在中心纤维 $X_0$ 上的限制 $E_0$ 数值有效, 则它在一般纤维 $X_\eta$ 上的限制 $E_\eta$ 也数值有效.
\end{proposition}
\begin{proof}
这由命题 \bookref{B.2.17} 推出.
\end{proof}
\begin{epigraph}
如果你不喜欢你的分析学家, 就去找身边的代数学家吧!
\epigraphauthor{Gert Almkvist}
\end{epigraph}
\par\clearpage
\endgroup
